% Front-matter prose, in the slots the UCL template defines.
% The template's Dedication slot is not used; this dissertation has none. Its
% Impact Statement slot is a PhD requirement, and for a taught Masters the
% department asks instead for a statement of what the work makes available, so
% that slot carries the data-and-code statement under its own heading.

\clearpage
{
    \pagestyle{abstract}
    \chapter*{Abstract}
    \addcontentsline{toc}{chapter}{Abstract}

A web agent looks at a web page, decides what to click or type, acts, and
repeats until a task is done. Before it can act it must be handed some
encoding of the page: cheap text listing elements and labels, an expensive
annotated screenshot, or both together. This choice is normally fixed once and
paid for at every step. This dissertation asks whether the expensive choice is
necessary at every step, and whether the steps that need it can be identified
cheaply enough to be worth identifying.

We hold the agent fixed and vary only what it is shown. Six \emph{observation
modes} differ in the page text supplied, the instruction wording, and whether
an annotated image is sent at all; task, actions, step limit and cost
accounting are identical across all six. We run all six on 8
website-and-model combinations, drawn from two public benchmarks
(VisualWebArena, WebArena) and three backbone models, and measure three
quantities in sequence: an \emph{oracle}'s per-task gain over the best single
fixed mode; whether that choice is predictable in advance from information a
deployed system would have; and whether acting on the prediction beats always
using the cheapest mode, once cost is counted.

The answers are, in order, \emph{yes}, \emph{no}, and \emph{the question is
mis-posed}. The oracle gains $+3.45$ to $+16.35$ percentage points of success
while spending 1.6--35.3\% less, in every combination on both benchmarks ---
though our own noise floor deflates this: rerunning one unchanged mode already
buys 2.0--7.6 points by itself, and two runs cannot pin down how much of that
ceiling repetition alone would reach, so the headroom's \emph{price} is the
safer claim, not its exact height. A policy counts as a win only if it is
\emph{both} more successful and cheaper than the free alternative --- a
\emph{Pareto} improvement. A learned router asked only whether \emph{any} mode
will solve a task, under a fully nested cross-validation protocol, clears that
bar in \textbf{0 of 8} combinations. Told the answer in advance instead of
predicting it, the same policy still clears the bar in only 1 of 8, recasting
the negative result from ``the learner failed'' to ``no Pareto-improving
\emph{deployable} operating point of this kind exists''. A stronger policy,
told \emph{which} mode solved each task rather than merely whether one did,
clears the bar in 7 of 8 --- but it needs an answer the benchmarks cannot
supply.

What blocks the prediction is training-example supply, not model capacity: at
2--27\% base success, 4 of 6 combinations lack enough solved tasks per mode to
train a classifier at all, and the strongest apparent signal on
VisualWebArena is the benchmark's own human-written difficulty annotation,
which no deployment has. Deliberately reduced training data confirms scarcity
is the mechanism and prices it: the benchmark that would relieve it needs to
be 2.1--4.2 times the size of the ones that exist --- a specification, not an
impossibility.

The contribution is not a router. It is a measured boundary: the conditions
under which \emph{representation routing}, meaning the choice of page encoding
per task, cannot be learned, with the
opportunity proven to exist before it is declared unreachable, and the failure
traced to a mechanism we can price.

    \thispagestyle{abstract}
}
\newpage

\clearpage
{
    \pagestyle{acknowledgements}
    \chapter*{Acknowledgements}
    \addcontentsline{toc}{chapter}{Acknowledgements}

I thank my supervisor, Mar\'ia P\'erez-Ortiz, for consistently redirecting
attention from what the system could be made to do towards what the evidence
could actually support, and my co-supervisor, Zekun Wu, for the structural
writing guidance that shaped how this document separates what was measured from
what was inferred.

I thank my parents, whose support made this degree possible, and You Wu, my
girlfriend, who lived alongside every deadline in it.

Compute for the local backbones was provided by the Holistic AI GB10 nodes and
by a dedicated A100 instance on UCL's ARC Condenser service; neither was shared
during the paper-grade runs, which matters for the latency and energy figures
reported in Chapter~\ref{ch:system}.

    \thispagestyle{acknowledgements}
}
\newpage

\clearpage
{
    \pagestyle{access}
    \chapter*{Data, Code and Public Access}
    \addcontentsline{toc}{chapter}{Data, Code and Public Access}

Every component of this study is publicly obtainable with one stated exception.
The benchmarks (VisualWebArena, WebArena) are open-source and self-hosted from
their published Docker images; two of the three backbones
(\texttt{Qwen3-VL-4B}, \texttt{google/gemma-3-4b-it}) are openly available model
weights run locally; the experiment harness, analysis pipeline and figure
scripts are released with the submission at
\url{https://github.com/Quarkgluonmixture/Cost-Aware-Routing-for-Web-Usage-Agents}.

\textbf{The exception is the large backbone.} \texttt{Qwen3-VL-235B-A22B} was
served through a metered API proxy hosted on AWS rather than self-hosted. Its outputs,
per-step logs and billing records are archived and released, but a third party
cannot re-run those conditions without equivalent proxy access. Results that
depend on that backbone alone are therefore \emph{replicable from our logs} but
not \emph{independently reproducible end to end}. Every claim that rests on
that backbone is marked where it is made, and
Chapters~\ref{ch:repr}--\ref{ch:e2e} are built so that none of the conclusions
depends on a single backbone.

    \thispagestyle{access}
}
\newpage
